% !TeX root = ../../../main.tex
\section{Conclusion}

The opportunity to use quantum computing becomes visible when a problem contains a mechanism that remains coherent, compact, and useful all the way through the machine. Periodicity can become phase. A rare event can become amplitude. A Hamiltonian can become controlled evolution. A symmetry can remove invalid states. A many-body quantum system can be represented without listing every amplitude. These are decisive abstractions because they change what must be computed, not merely where it is stored.

The negative criteria are equally important. Large data are not enough. Generic hardness is not enough. A QUBO is not enough. Entanglement is not enough. A circuit that fits is not enough. The data must enter, the oracle must run, garbage must be uncomputed, the observable must be measured, the errors must fit a budget, and the answer must beat a real classical opponent at a useful scale.

Subatomic physics offers a natural home for these ideas because the underlying dynamics are already quantum, although continuum limits, state preparation, and fault-tolerant resources remain formidable. Astronomy offers narrower but fascinating paths: structured matched filtering, quantum kernels inside multiscale physical models, and quantum-enhanced sensing before light becomes a classical file. In each case, the right question is not ``Where can I use a quantum computer?'' It is ``Where is the coherent bottleneck, what compact truth can it reveal, and what does the entire experiment cost?''

Planck teaches us to respect the anomaly. Einstein teaches us to design the killing thought experiment. Heisenberg teaches us to begin with observables. Schr\"odinger teaches us to choose the representation that exposes the spectrum. Born teaches us to calculate uncertainty. Pauli teaches us that constraints define the valid world. Feynman teaches us to open the drawer.

That last lesson may be the most practical. A quantum algorithm should survive an engineer reaching around the back of it.
